\subsection{Baselines and controls}\label{sec:met-baselines}
\textbf{Fixed-time} is the floor every deployed junction achieves: each
network's own default cyclic plan, the same phase menu as every other arm.
\textbf{MaxPressure} is the theoretically-grounded adaptive baseline
\cite{varaiya2013maxpressure}, and doubles as the shield's fallback policy.
The control for every fine-tuning comparison is the \emph{same base model
compiled through the identical int4 pipeline} (denoted stock, or ft0), not
the vendor catalogue artifact, for two reasons. I observed the catalogue
Qwen3-0.6B GPU artifact emitting corrupted tokens after long service
uptime, an unlogged operational observation rather than a measured finding
(\S\ref{sec:res-reliability}). Independently, and sufficient on its own,
that artifact differs from my compiled one in quantisation recipe even
when healthy. Holding the compilation pipeline fixed leaves the fine-tuned
weights as the single manipulated variable.

One model class is excluded by measurement rather than swept:
phi-4-mini-reasoning's p99 decision latency (11.1\,s) exceeds the 10\,s
decision interval on this hardware, so no deployment could serve it in
time. That is a deployability finding, not a control one, since the
harness blocks on the model call with simulated time stopped and a
closed-loop cell would flatter a slow model rather than penalise it; the
p99 is one format of four on one seed, the other three at 8.69, 9.61 and
9.74\,s, and the model is recorded as latency-excluded with its single
corridor probe retained descriptively. Two further controls target specific
hypotheses: a \textbf{leave-one-topology-out} (LOTO) student, trained with
all Old Street states removed, tests whether distillation gains are
topology memorisation, and four \textbf{format-specialist} students, each
trained on one prompt format alone, test whether a deployed junction would
prefer a per-format model over a generalist. The answer
(\S\ref{sec:res-offline}) was no, at the single scale (3.8B) that sweep
ran at.
